\section{Partons as quanta in infinite-momentum states and
large-momentum expansion}
\label{sec:imf}

Let us begin with a brief overview of the parton formalism. In the textbooks, PDFs are usually defined in terms of
LF correlations in QCD~\cite{Sterman:1994ce,Collins:2011zzd}.
The LF is defined by $t-z=$~constant, if a massless particle
is travelling along the $z$-direction, with variations of other
coordinates, $t+z$ and transverse-space dimensions, defining
a three-dimensional front surface.
Introduce two independent LF four-vectors with dimension-one
parameter $\Lambda$,
\begin{eqnarray}
   p^\mu &=& (\Lambda, 0, 0, \Lambda )\ , \nonumber \\
   n^\mu &=& \left(1/(2\Lambda), 0, 0, -1/(2\Lambda)\right)\ ,
\end{eqnarray}
then $p^2=n^2=0$, and $p\cdot n=1$. Different LFs
are defined by different coordinate distance $\lambda$
along the $n$-direction.

Consider now the quark PDFs in a state
$|P\rangle$ with mass $M$ and four-momentum $P^\mu= (P^0,0,0,P^z)
= p^\mu + (M^2/2) n^\mu$, which can be used to solve for $\Lambda =(P^0+P^z)/2$.
Using $\psi$ to denote a full-QCD quark field,
the LF correlation function in coordinate space is,
\begin{equation}
     h(\lambda)= \frac{1}{2P^+}\langle P|\overline \psi(\lambda n)\Gamma W(\lambda n, 0)\psi(0)|P\rangle \ ,
\label{Eq:correlationm}
\end{equation}
where $\Gamma$ is a Dirac matrix, $W$ is a straight Wilson-line gauge link,
and $\lambda$ is the LF distance. All other coordinates have been taken
to be zero. Due to the invariance of the LF under Lorentz boosts along the $z$-direction, the
above correlation function is independent of the residual momentum $P^z$. Quite often, $P^z$ is taken to be zero.

The quark PDF is just the Fourier transform of the above LF correlation~\cite{Collins:2011zzd},
\begin{equation}
     f(x) =\int^\infty_{-\infty} \frac{d\lambda}{2\pi} e^{-ix\lambda}h(\lambda)  \ .
\end{equation}
In this way, partons can be studied without using the EFT machinery although they are effective degrees of freedom (dof's)
to describe the LF collinear modes. The reason is that, the parton dof's are automatically projected out through the LF correlators applied to the full QCD state $|P\rangle$.
On the other hand, these parton dof's can also be explicitly separated in the QCD Lagrangian, as is done in soft-collinear effective theory (SCET) where they are represented by LF collinear
fields~\cite{Bauer:2000yr,Bauer:2001ct,Bauer:2001yt}.

In the traditional parton formalism, the correlations are time-dependent, or in other words, the operators are in the Heisenberg picture. As such, we say that the formalism
is Minkowskian and thus difficult for Monte Carlo simulations due to
the famous ``sign'' problem. If one chooses
$\xi^- = (t+z)/\sqrt{2}$ as the ``new time'' coordinate,
and integrates it out, one obtains a Hamiltonian formalism
for partons, which has been called LF quantization (LFQ) in the
literature~\cite{Dirac:1949cp}. LFQ
is also a very difficult formalism to work with, despite the fact that
much progress has been made~\cite{Brodsky:1997de}.

An alternative parton formalism can be obtained
by adapting Feynman's original idea about partons
to the context of a field theory~\cite{Ji:2014gla,Ji:2020ect}.
Feynman considered~\cite{Feynman:1973xc} the momentum distribution of a composite system,
$f(k^z, P^z)$, where $P^z$ is the center-of-mass momentum and $k^z$ the longitudinal momentum carried by the parton whose
transverse momentum $\vec{k}_\perp$ has been integrated over.
The $P^z$-dependence of the momentum distribution is clearly
a relativistic effect: According to Poincar\'e symmetry,
the Hamiltonian of a system depends on the frame, and changes under Lorentz boosts
according to,
\begin{equation}
  [H, K^i] = i P^i\ ,
\end{equation}
where $K^i$ $(i = 1,2,3)$ are the boost operators. Therefore,
the wave functions are frame-dependent, leading to frame-dependent
momentum distributions. Because $K^i$ depends on interactions, the
frame-dependence is a dynamical problem, and generally requires
non-perturbative solutions.

An important feature of the momentum distribution of a system is that it is a static or time-independent quantity. In QCD,
it is related to the following spatial correlation,
\begin{equation}
    \tilde h(z,P^z)= \frac{1}{2N}\langle P|\overline \psi(z)\Gamma W(z, 0)\psi(0)|P\rangle \ ,
\label{eq:esme}
\end{equation}
where $W(z,0)$ is a spacelike, straight-line gauge link, and $N$ is normalization factor depending on the Dirac matrix.
Feynman then considered the infinite-momentum limit, assuming that such a limit
exists,
\begin{equation}
  P^z \to \infty, ~~~~~z\to 0, ~~~~~\lambda = zP^z ~{\rm finite}\ ,
\end{equation}
i.e., the relevant correlation function for partons is
\begin{equation}
  \tilde h(\lambda= zP^z) = \langle P^z=\infty|\overline \psi(z)\Gamma W(z, 0)\psi(0)|P^z=\infty\rangle\ .
\label{Eq:correlatione}
\end{equation}
It is clear that in field theories this is a non-trivial limit. In fact,
it can be shown that such a limit only exists in asymptotically-free theories, where
the high-momentum modes are perturbative~\cite{Ji:2020ect}.

If one ignores the subtlety of the limit,
the correlation in Eq.~(\ref{Eq:correlatione})
is related to that  in Eq.~(\ref{Eq:correlationm})
by an infinite Lorentz transformation~\cite{Ji:2013dva}.
Our ``new'' form of parton formalism works with time-independent
correlators and the IMF wave function.
Since the operator is time-independent, it is the
Schr\"odinger representation of parton physics
if an analogy between time-translation and Lorentz boost is
made~\cite{Ji:2020ect}.

In QCD, however, the correlations $\tilde h(\lambda)$
and $h(\lambda)$ are different. The difference arises
from the presence of the UV cut-off. In the physical
momentum distribution, the cut-off must always be much
larger than the hadron momentum. As a result, the parton momentum
is allowed to be larger than the hadron momentum, or $|x|$ can be
larger than 1, without
violating any laws of physics. On the other hand, the standard PDFs
have support $|x|\le1$, corresponding to a UV cut-off smaller than
the hadron momentum. Thanks
to the asymptotic freedom, these two different UV limits
can be connected to each other by perturbation theory in QCD. This makes it possible to extract LF parton physics defined in Eq. (\ref{Eq:correlationm}) from the Euclidean form in Eq. (\ref{Eq:correlatione}).

Eq. (\ref{Eq:correlatione}) is the starting point
of the LaMET expansion, where we first
compute the quasi-LF correlation functions at a finite,
but large momentum $P^z\gg\Lambda_{\rm QCD}$.
To make the expansion work, in principle one needs ${\tilde h}(z, P^z)$ with $-\infty<z<\infty$, or
${\tilde h}(\tilde \lambda, P^z)$ at all quasi-LF distances $\tilde \lambda=z P^z$.
While in reality, of course, due to
the finite volume, lattice data will always stop at some
large $z(\tilde \lambda)$ which we call $z_{\rm L}(\tilde \lambda_{\rm L})$. We will
deal with issues of finite $\tilde \lambda_{\rm L}$ later.
For the discussion in this section, we assume that $\tilde h(\tilde \lambda, P^z)$
is known in $[-\infty,\infty]$, i.e., in the whole $\tilde \lambda$ range at a large
$P^z$.

With the above quasi-LF correlation, one can make a straightforward
Fourier transformation
\begin{equation}
    {\tilde f}(y,P^z) =\int^\infty_{-\infty} \frac{d\tilde \lambda}{2\pi}\, e^{i\tilde \lambda y}\, \tilde h(z, P^z)\ .
\end{equation}
The physical interpretation of ${\tilde f}(y,P^z)$ hinges on the large momentum
expansion~\cite{Ji:2013dva,Xiong:2013bka,Ma:2014jla,Izubuchi:2018srq}
\begin{align}\label{matching}
    {\tilde f}(y, P^z) &= \int^1_{-1} dx\, C\Big(\frac y x, \frac{xP^z}{\mu}\Big) f(x, \mu) \nonumber \\
    &+ {\cal O}\Big(\frac{\Lambda^2_{\rm QCD}}{y^2(P^z)^2}, \frac{\Lambda^2_{\rm QCD}}{(1-y)^2(P^z)^2}\Big),
\end{align}
where $\mu$ is a factorization scale, $\Lambda_{\rm QCD}$ is the hadronic scale, and
$f(x,\mu)$ is the standard PDF that can be extracted from the above equation.
The large scales $(yP^z)^2$ and $((1-y)P^z)^2$ are associated with the active quark
and the spectator momenta, respectively. According to the standard EFT methodology,
any large scale that is not forbidden shall be allowed in the expansion, and the linear
dependence is absent in dimensional regularization due to space-time symmetry.
Therefore, the validity of this expansion relies on the smallness of the expansion
parameters $\Lambda^2_{\rm QCD}/[y^2(P^z)^2]$ and $\Lambda^2_{\rm QCD}/[(1-y)^2(P^z)^2]$.

The $C$ factor in the above equation can be calculated perturbatively. At leading-order in $\alpha_s$, we can identify $\tilde f(y, P^z)$ with $f(y,\mu)$ (ignoring the power corrections for the moment), thus they have the same asymptotic behavior as $y\to 0$ and $y\to1$. Beyond leading-order, this will be changed by perturbative corrections. To see this, let us take the following simple form of $f(x, \mu)$ as an example
\begin{equation}
x^a (1-x)^b,
\end{equation}
with $a, b$ controlling the asymptotic behavior at $x\to 0$ and $x\to 1$, respectively. The perturbative one-loop corrections lead to the following change in the asymptotic behavior
\begin{align}
\delta\tilde f(y, P^z)\sim \alpha_s y^a \ln y \ \  {\rm as}\ \  y\to 0\,.
\end{align}
When resummed to all orders in perturbation theory, this yields a power law behavior of the form $\tilde f(y, P^z)\sim  y^{a+\gamma}$ with $\gamma$ being associated with the anomalous dimension of the operator defining $\tilde f(y, P^z)$. Similar behavior also occurs as $y\to 1$ for realistic PDFs with $b>0$. This can also be seen from the coordinate space analysis to be presented below.

Therefore, for a given large $P^z$,
there is a range of $y$ where high-order corrections as well as power corrections are small, and this range can be translated into
a valid range $x$ for the PDFs. Thus, one can systematically obtain the PDFs in an interval $[x_{\rm min}, x_{\rm max}]$ ($x_{\rm min}$ will approach 0 and $x_{\rm max}$ approach 1 as $P^z\to \infty$).
In other words, the LaMET
expansion provides a natural way to calculate parton distributions
in an interval of the parton momentum $x$, similar to extracting parton distributions from
experimental data at finite energies.
